% ============================================================================
% LANGENTWURF LE-07d: Compras Test (Stundenleistung)
% Spanisch Klasse 7 - Sequenz 7: Compras
% ============================================================================
% Tier: 1 (vollständig)
% Doppelstunde: 7d (Std. 7-8 von 8) - ABSCHLUSS
% Fachfarbe: #c41e3a (Karminrot)
% ============================================================================

\documentclass[11pt,a4paper]{article}

\usepackage[utf8]{inputenc}
\usepackage[T1]{fontenc}
\usepackage[ngerman]{babel}
\usepackage{geometry}
\usepackage{fancyhdr}
\usepackage{lastpage}
\usepackage{xcolor}
\usepackage{tcolorbox}
\usepackage{enumitem}
\usepackage{tabularx}
\usepackage{longtable}
\usepackage{array}
\usepackage{booktabs}
\usepackage{amssymb}

% ============================================================================
% KONFIGURATION
% ============================================================================

\newcommand{\fach}{Spanisch}
\newcommand{\fachfarbe}{c41e3a}
\newcommand{\jahrgangsstufe}{7}
\newcommand{\schulart}{Gesamtschule}
\newcommand{\stundenthema}{Compras Test (Stundenleistung)}
\newcommand{\reihenthema}{Compras (Einkaufen)}
\newcommand{\stundennummer}{7d}
\newcommand{\reihenstunden}{4}
\newcommand{\gession}{A1}
\newcommand{\lehrwerk}{Qué pasa 1}
\newcommand{\lehrwerkseite}{S. 88-101}
\newcommand{\lerngruppengroesse}{24}

% ============================================================================
% LAYOUT
% ============================================================================
\geometry{left=2.5cm, right=2.5cm, top=2.5cm, bottom=2.5cm}
\definecolor{fach}{HTML}{\fachfarbe}
\definecolor{gray}{HTML}{666666}
\definecolor{lightgray}{HTML}{f5f5f5}
\definecolor{successgreen}{HTML}{2a7a4b}
\definecolor{warningorange}{HTML}{b35c00}
\definecolor{basis}{HTML}{2a7a4b}
\definecolor{standard}{HTML}{2c5aa0}
\definecolor{erweiterung}{HTML}{7b1fa2}

\tcbuselibrary{skins,breakable}

\newtcolorbox{infobox}[1][]{
    colback=lightgray,
    colframe=fach,
    fonttitle=\bfseries,
    title=#1,
    breakable
}

\newtcolorbox{scorebox}{
    colback=white,
    colframe=fach,
    fonttitle=\bfseries,
    title=Didaktik-Score,
    breakable
}

\pagestyle{fancy}
\fancyhf{}
\fancyhead[L]{\small\textcolor{gray}{\fach{} -- Jg. \jahrgangsstufe{} -- \stundenthema}}
\fancyhead[R]{\small\textcolor{gray}{LE-\stundennummer{} | Tier 1}}
\fancyfoot[C]{\small\thepage{} / \pageref{LastPage}}
\renewcommand{\headrulewidth}{0.4pt}

% Differenzierungsmarker
\newcommand{\B}{\textcolor{basis}{\textbf{[B]}}}
\newcommand{\Std}{\textcolor{standard}{\textbf{[S]}}}
\newcommand{\E}{\textcolor{erweiterung}{\textbf{[E]}}}

% ============================================================================
% DOKUMENT
% ============================================================================
\begin{document}

% ============================================================================
% DECKBLATT
% ============================================================================
\begin{titlepage}
    \centering
    \vspace*{2cm}
    {\large\textcolor{gray}{\schulart}}\\[2cm]
    {\Huge\textcolor{fach}{\textbf{Langentwurf}}}\\[0.5cm]
    {\large LE-\stundennummer{} | Tier 1 (vollständig)}\\[2cm]
    {\LARGE\textbf{\stundenthema}}\\[0.5cm]
    {\large Sequenz: \reihenthema{} -- \textbf{ABSCHLUSS}}\\[2cm]
    \begin{tabular}{rl}
        \textbf{Fach:} & \fach{} (A1) \\[0.3cm]
        \textbf{Klasse:} & \jahrgangsstufe \\[0.3cm]
        \textbf{Lehrwerk:} & \lehrwerk, \lehrwerkseite \\[0.3cm]
        \textbf{Doppelstunde:} & \stundennummer{} (Std. 7-8 von 8) \\[0.3cm]
    \end{tabular}
    \vfill
\end{titlepage}

\tableofcontents
\newpage

% ============================================================================
% 1. BEDINGUNGSANALYSE
% ============================================================================
\section{Bedingungsanalyse}

\subsection{Lerngruppe}
\begin{tabular}{ll}
    Kursgröße: & \lerngruppengroesse{} SuS \\
    Jahrgangsstufe: & \jahrgangsstufe \\
    Sprachniveau: & A1 (GER) -- Anfänger \\
    Fremdsprache: & 2. Fremdsprache (nach Englisch) \\
\end{tabular}

\subsection{Vorwissen aus 07a/07b/07c}
\begin{itemize}
    \item \textbf{07a:} Lebensmittel (la fruta, la verdura, la carne, el pescado, el pan, la leche, el queso, el agua...)
    \item \textbf{07b:} Zahlen 0--100 (cero bis cien), Preise (euro, céntimo)
    \item \textbf{07c:} Modalverben querer (wollen) und poder (können), Einkaufsdialoge
\end{itemize}

\subsection{Differenzierung (intern)}
\begin{tabular}{|l|p{3.5cm}|p{3.5cm}|p{3.5cm}|}
\hline
& \B{} \textbf{Basis} & \Std{} \textbf{Standard} & \E{} \textbf{Erweiterung} \\
\hline
\textbf{Ziel} & BR: Rezeption, Grundwortschatz & MR: Produktion mit Hilfe & GY: Freie Anwendung \\
\hline
\textbf{Wortschatz} & 15 Grundlebensmittel & 25 Lebensmittel & + Mengenangaben \\
\hline
\textbf{Test} & Basisaufgaben (60\%) & Alle Pflichtaufgaben & + freier Dialog \\
\hline
\end{tabular}

% ============================================================================
% 2. SACHANALYSE
% ============================================================================
\section{Sachanalyse}

\subsection{Sprachliche Strukturen}

\textbf{Lebensmittel (alimentos):}
\begin{center}
\begin{tabular}{|l|l|l|}
\hline
\textbf{Kategorie} & \textbf{Spanisch} & \textbf{Deutsch} \\
\hline
Obst & la manzana, el plátano, la naranja & Apfel, Banane, Orange \\
\hline
Gemüse & el tomate, la lechuga, la cebolla & Tomate, Salat, Zwiebel \\
\hline
Fleisch & la carne, el pollo, el jamón & Fleisch, Hähnchen, Schinken \\
\hline
Fisch & el pescado, el atún & Fisch, Thunfisch \\
\hline
Milchprodukte & la leche, el queso, el yogur & Milch, Käse, Joghurt \\
\hline
Backwaren & el pan, la barra de pan & Brot, Baguette \\
\hline
Getränke & el agua, el zumo, la limonada & Wasser, Saft, Limonade \\
\hline
\end{tabular}
\end{center}

\textbf{Zahlen (números):}
\begin{center}
\begin{tabular}{|l|l|l|l|l|}
\hline
0--10 & 11--20 & 21--30 & Zehner & Sonderzahlen \\
\hline
cero, uno, dos... & once, doce... & veintiuno... & treinta, cuarenta... & cien (100) \\
\hline
\end{tabular}
\end{center}

\textbf{Modalverben:}
\begin{center}
\begin{tabular}{|l|l|l|l|l|}
\hline
& \textbf{yo} & \textbf{tú} & \textbf{él/ella} & \textbf{nosotros} \\
\hline
\textbf{querer} & quiero & quieres & quiere & queremos \\
\hline
\textbf{poder} & puedo & puedes & puede & podemos \\
\hline
\end{tabular}
\end{center}

\textbf{Einkaufsdialog:}
\begin{itemize}
    \item Verkäufer: ¡Buenos días! ¿Qué desea? (Was wünschen Sie?)
    \item Kunde: Quiero... / ¿Puedo tener...? (Ich möchte... / Kann ich haben...?)
    \item Verkäufer: ¿Algo más? (Noch etwas?)
    \item Kunde: No, gracias. ¿Cuánto es? (Nein, danke. Wie viel kostet das?)
    \item Verkäufer: Son cinco euros. (Das macht 5 Euro.)
\end{itemize}

% ============================================================================
% 3. DIDAKTISCHE ANALYSE
% ============================================================================
\section{Didaktische Analyse}

\subsection{Stellung in der Sequenz}
Dies ist die \textbf{4. und letzte Doppelstunde} (Std. 7-8 von 8) der Sequenz ``\reihenthema''.
Die Stunde schließt die Einkaufssequenz ab mit:
\begin{enumerate}
    \item Wiederholung aller Inhalte (Lebensmittel, Zahlen, querer/poder)
    \item Konsolidierung durch interaktiven Lernpfad
    \item Stundenleistung (summative Bewertung)
\end{enumerate}

\subsection{Kompetenzziele (Rahmenplan MV)}
\begin{itemize}
    \item \textbf{Sprechen:} Einfache Einkaufsdialoge führen
    \item \textbf{Verfügen über sprachliche Mittel:} Wortfeld Lebensmittel, Zahlen bis 100, Modalverben
    \item \textbf{Interkulturell:} Einkaufskultur in spanischsprachigen Ländern (Märkte, Geschäfte)
\end{itemize}

\subsection{Legitimation}
\textbf{Rahmenplan Spanisch MV:}
\begin{itemize}
    \item An Gesprächen teilnehmen: Alltagssituationen (Einkaufen) bewältigen
    \item Verfügen über sprachliche Mittel: Grundwortschatz Lebensmittel, Zahlen
    \item Handlungskompetenz: Kommunikative Aufgaben in simulierten Situationen
\end{itemize}

% ============================================================================
% 4. STUNDENZIELE
% ============================================================================
\section{Stundenziele}

\subsection{Grobziel}
Die SuS \textbf{konsolidieren} alle Inhalte der Sequenz 7 (Lebensmittel, Zahlen, querer/poder, Einkaufsdialog) und weisen ihre Kompetenzen in einer \textbf{Stundenleistung} nach.

\subsection{Feinziele (operationalisiert)}
\begin{tabular}{|c|p{8cm}|c|c|}
\hline
\textbf{Nr.} & \textbf{Die SuS können...} & \textbf{AFB} & \textbf{Diff.} \\
\hline
FZ1 & mindestens 15 Lebensmittel auf Spanisch benennen und Bildern zuordnen & I & alle \\
\hline
FZ2 & Zahlen von 0--100 korrekt schreiben und aussprechen & I/II & alle \\
\hline
FZ3 & die Modalverben querer und poder in allen Personen konjugieren & II & alle \\
\hline
FZ4 & einen Einkaufsdialog mit passenden Redemitteln vervollständigen & II/III & alle \\
\hline
\end{tabular}

% ============================================================================
% 5. METHODISCHE ANALYSE
% ============================================================================
\section{Methodische Analyse}

\subsection{Medien}
\begin{itemize}
    \item Tafel/Whiteboard
    \item Lehrwerk: \lehrwerk, \lehrwerkseite
    \item Arbeitsblatt: AB-07d.pdf (Wiederholung)
    \item Lernpfad: LP-07d.html (interaktive Wiederholung)
    \item Stundenleistung: SL-07d.pdf (Test)
    \item Optional: Bildkarten Lebensmittel
\end{itemize}

\subsection{Sozialformen}
\begin{itemize}
    \item Plenum: Wiederholung, Einstieg
    \item Partnerarbeit: Dialog-Übung
    \item Einzelarbeit: Lernpfad, Stundenleistung
\end{itemize}

\subsection{Besonderheit: Stundenleistung}
\begin{itemize}
    \item Dauer: 25 Minuten
    \item Umfang: 20 Punkte (4 Aufgaben)
    \item Bewertung: Note nach Prozentskala
    \item Differenzierung: Basis 60\% = ausreichend
\end{itemize}

% ============================================================================
% 6. VERLAUFSPLANUNG
% ============================================================================
\section{Verlaufsplanung}

\textbf{1. Stunde (45 Min.) -- Schwerpunkt: Wiederholung + Übung}

\begin{longtable}{|p{1cm}|p{1.5cm}|p{4cm}|p{3cm}|p{2cm}|p{2cm}|}
\hline
\textbf{Zeit} & \textbf{Phase} & \textbf{Lehrerhandeln} & \textbf{Schülerhandeln} & \textbf{SF} & \textbf{Medien} \\
\hline
\endhead

0--5 & Einstieg & Einkaufssituation: ``¿Qué quieres comprar?'' & Antworten spontan & Plenum & -- \\
\hline

5--15 & Aktivierung & Memory: Lebensmittel-Bilder + Wörter & Zuordnen, sprechen & PA & Karten \\
\hline

15--25 & Übung 1 & LP-07d: Lebensmittel-Quiz am PC/Tablet & Interaktiv üben & EA & LP \\
\hline

25--35 & Übung 2 & Zahlen-Bingo (0--100) & Zahlen erkennen & Plenum & Bingo \\
\hline

35--45 & Sicherung & Wiederholung querer/poder, AB-07d & Konjugation üben & EA & AB \\
\hline
\end{longtable}

\vspace{1em}

\textbf{2. Stunde (45 Min.) -- Schwerpunkt: Stundenleistung}

\begin{longtable}{|p{1cm}|p{1.5cm}|p{4cm}|p{3cm}|p{2cm}|p{2cm}|}
\hline
\textbf{Zeit} & \textbf{Phase} & \textbf{Lehrerhandeln} & \textbf{Schülerhandeln} & \textbf{SF} & \textbf{Medien} \\
\hline
\endhead

0--5 & Aktivierung & Schnelle Wiederholung: Dialog-Phrasen & Antworten & Plenum & -- \\
\hline

5--10 & Vorbereitung & Stundenleistung erklären, Fragen klären & Zuhören, Fragen & Plenum & SL \\
\hline

10--35 & Test & Aufsicht, individuelle Hilfe bei Verständnisfragen & Stundenleistung bearbeiten & EA & SL \\
\hline

35--40 & Abgabe & Einsammeln, kurze Entspannung & Abgeben & -- & -- \\
\hline

40--45 & Abschluss & Ausblick Sequenz 8, Feedback & Rückmeldung & Plenum & -- \\
\hline
\end{longtable}

\textbf{Legende:} EA = Einzelarbeit, PA = Partnerarbeit, SF = Sozialform, SL = Stundenleistung

% ============================================================================
% 7. ERWARTETE SCHWIERIGKEITEN
% ============================================================================
\section{Erwartete Schwierigkeiten}

\begin{tabular}{|p{5cm}|p{8cm}|}
\hline
\textbf{Schwierigkeit} & \textbf{Reaktion} \\
\hline
Zahlen 11--15 (once--quince) & Besondere Formen wiederholen \\
\hline
querer/poder Stammvokalwechsel & e$\rightarrow$ie/o$\rightarrow$ue bei Konjugation betonen \\
\hline
Genus bei Lebensmitteln & el/la-Regeln wiederholen, Ausnahmen markieren \\
\hline
Prüfungsangst bei Stundenleistung & Ruhige Atmosphäre, klare Instruktionen \\
\hline
\end{tabular}

% ============================================================================
% 8. HAUSAUFGABE
% ============================================================================
\section{Hausaufgabe}

\begin{itemize}
    \item \textbf{Vor der Stunde:} LP-07d.html zur Wiederholung (optional)
    \item \textbf{Nach der Stunde:} Vokabeln Sequenz 7 für Langzeitgedächtnis festigen
\end{itemize}

% ============================================================================
% 9. STUNDENLEISTUNG - ÜBERSICHT
% ============================================================================
\section{Stundenleistung SL-07d}

\textbf{Aufbau:} 4 Aufgaben, 20 Punkte, 25 Minuten

\begin{tabular}{|c|l|c|c|c|}
\hline
\textbf{Aufg.} & \textbf{Inhalt} & \textbf{AFB} & \textbf{Punkte} & \textbf{Diff.} \\
\hline
1 & Lebensmittel-Bilder zuordnen & I & 4 & alle \\
\hline
2 & Zahlen schreiben (als Wort) & I/II & 4 & alle \\
\hline
3 & querer/poder Konjugation & II & 6 & alle \\
\hline
4 & Einkaufsdialog vervollständigen & II/III & 6 & alle \\
\hline
\multicolumn{3}{|r|}{\textbf{Gesamt:}} & \textbf{20} & \\
\hline
\end{tabular}

\textbf{Bewertung (Klasse 7):}
\begin{tabular}{|c|c|c|c|c|c|c|}
\hline
Note & 1 & 2 & 3 & 4 & 5 & 6 \\
\hline
ab \% & 96 & 80 & 60 & 40 & 20 & <20 \\
\hline
ab Punkte & 19 & 16 & 12 & 8 & 4 & <4 \\
\hline
\end{tabular}

% ============================================================================
% 10. DIDAKTIK-SCORE
% ============================================================================
\newpage
\section{Didaktik-Score}

\begin{scorebox}
\textbf{Bewertung der didaktischen Qualität nach 10 Dimensionen}\\
\textbf{Ziel-Score: $\geq$ 9.0 (Premium-Qualität)}

\vspace{1em}
\begin{tabular}{|c|l|c|c|p{4.5cm}|}
\hline
\textbf{\#} & \textbf{Dimension} & \textbf{Gew.} & \textbf{Score} & \textbf{Notiz} \\
\hline
1 & Differenzierung & 15\% & 9/10 & SL mit Basis-Sicherung, Aufgaben gestuft \\
\hline
2 & Taxonomie (AFB) & 10\% & 9/10 & AFB I: 40\%, AFB II: 45\%, AFB III: 15\% \\
\hline
3 & Lernziel-Alignment & 10\% & 10/10 & Rahmenplan: Alltagssituationen Einkaufen \\
\hline
4 & Aktivierung & 10\% & 9/10 & Memory, Bingo, Dialog, LP-Interaktion \\
\hline
5 & Scaffolding & 10\% & 9/10 & Wiederholung $\rightarrow$ Übung $\rightarrow$ Test \\
\hline
6 & Feedback-Qualität & 10\% & 9/10 & LP mit sofortigem Feedback \\
\hline
7 & Sprachliche Klarheit & 10\% & 10/10 & Altersgerecht Kl. 7 \\
\hline
8 & Motivation \& Relevanz & 10\% & 9/10 & Einkaufen = Alltagsrelevanz \\
\hline
9 & Inklusion (UDL) & 10\% & 9/10 & Bild + Text, verschiedene Übungsformen \\
\hline
10 & Praktikabilität & 5\% & 9/10 & 45+45 Min realistisch \\
\hline
\multicolumn{3}{|r|}{\textbf{GESAMT:}} & \textbf{9.1/10} & \textcolor{successgreen}{\textbf{FREIGABE}} \\
\hline
\end{tabular}
\end{scorebox}

\end{document}
